% ============================================================================
\section{Experimental Setup}
\label{sec:experiments}
% ============================================================================

\subsection{A four-domain distillation setting}

We study mathematics, code, instruction following, and science with a
Qwen3-1.7B student in non-thinking mode. Mathematics and instruction
following use domain-specific Qwen3-1.7B RL teachers. Code and science
are routed to the same frozen Qwen3-4B model through separate task
endpoints. This setting combines specialized feedback and a larger
generalist teacher in one student.

The default loss is Equation~\ref{eq:dense_topk_loss} with teacher
$k=64$, and each task has weight $w_i=1/4$. Each step processes
$G=16$ micro-batches of $b=4$ responses, with $2\leq m_i\leq8$.
Uniform uses four micro-batches per task. Runs last $500$ steps;
each task supplies a shuffled stream of $16{,}000$ unique prompts,
and each response is capped at $4{,}096$ tokens. There is one optimizer
update per fresh rollout batch. Noise and time estimates use EMA decay
$0.9$. All runs use AdamW, the same task prompt orders, and one training
seed per method.

One $96$-GB GPU trains the student, and one $48$-GB GPU serves
student rollouts and teacher scoring. The two generalist task endpoints
share their teacher weights in memory. Appendix~\ref{app:reference_tables}
records the configuration.
\missing{Complete the exact checkpoint and data revisions, optimizer
settings, software versions, and measured resource use.}

\subsection{Comparisons}

The core comparison keeps the loss, response averaging, and task weights
fixed, varying how the $16$ micro-batches are assigned:
\begin{itemize}
    \item \textbf{Uniform}: four micro-batches per task.
    \item \textbf{GPAS (per step)}: counts from smoothed preconditioned
    gradient noise.
    \item \textbf{GPAS (cost-aware)}: counts from the time--variance
    criterion in Equation~\ref{eq:batch_cost_objective}.
    \item \textbf{Unpreconditioned allocation}: the same variance rule
    with $D=I$ for count selection; AdamW still performs the update.
    \item \textbf{D$^3$ signal, fixed weights}: counts from D$^3$-MOPD's
    remaining-gap and descent-velocity signal, combined with our fixed
    task-mean weights. This compares learning progress with gradient
    variability as a scheduling signal.
\end{itemize}

We also include two recipe-level comparisons.
\textbf{D$^3$-MOPD scheduling} uses the same gap--velocity scheduling rule but
averages responses across the step, so a task's effective weight follows
$m_i/G$. \textbf{Open-MOPD ($K=1$)} uses its student top-$16$ dense
loss, token-share balancing, and gap-aware weighting with one update per
rollout batch. Student reward refresh has no intervening update to
correct in this $K=1$ setting. These are adaptations to our model, data,
and training budget; their task objectives differ from the fixed-weight
comparison. Appendix~\ref{app:baselines} specifies the adaptations.

Finally, \textbf{TA-OPD + Uniform} and \textbf{TA-OPD + GPAS} use
the tail-aware loss of \citet{huang2026tailaware} at the same $k=64$.
This pair evaluates whether the allocation remains useful when the
dense supervision also accounts for the omitted probability mass.
It brings the study to nine training configurations.

\subsection{Capability, efficiency, and allocation traces}

\paragraph{Capability and learning curves.}
We report MATH-500 greedy pass@1, a fixed LiveCodeBench pass@1 slice,
IFBench strict accuracy, and GPQA-Diamond average@4. The initial student
and assigned teachers are evaluated with the same benchmark protocols.
The main table reports all four scores and their unweighted mean.
Every $50$ steps, we also evaluate $128$ fixed held-out prompts per task
using the common teacher top-$64$ loss and equal task weights.
This provides a common distillation metric even for recipes trained
with different losses or task weights. Shared evaluation seeds allow
paired prompt-level bootstrap intervals; these describe evaluation
uncertainty, not training-seed variation.
\missing{Insert benchmark versions, the code evaluation date range,
and complete decoding settings.}

\paragraph{Measured efficiency.}
We plot the common held-out loss against both optimizer steps and
GPU hours. GPU hours include all time on the two reserved devices,
including waiting, synchronization, statistics, and checkpointing.
For the fixed-weight default-loss methods, we also report the time to
Uniform's final held-out loss, interpolating between adjacent evaluations
and marking unreached targets. Capability scores accompany this
loss-based comparison because teacher agreement and task performance
need not improve together.

\paragraph{Understanding the allocation.}
The ordinary training loop logs task loss, its recent rate of change,
response length, preconditioned and unpreconditioned noise, assigned
counts, and time per micro-batch. These traces show when a task receives
more computation and whether the decision follows remaining gap,
gradient variability, or cost. We report the time and memory used by
the GPAS statistics as part of the implementation cost.

% ============================================================================
\section{Results}
\label{sec:results}
% ============================================================================

\subsection{Training efficiency and per-domain capability}

\begin{table}[t]
\centering
\small
\setlength{\tabcolsep}{3pt}
\caption{Capability and training efficiency in the four-domain study.
$F_{\rm eval}$ is the common equal-weight teacher top-$64$ evaluation
loss. Time to target is reported within the default-loss, fixed-weight
group. Dashes denote inapplicable entries; blank cells await measurements.}
\label{tab:main_results}
\begin{tabular}{lccccccc}
\toprule
Method & Math & Code & IF & Science & Mean & $F_{\rm eval}$ & GPU h to target \\
\midrule
Initial student & & & & & & & --- \\
Assigned teachers & & & & & & --- & --- \\
\midrule
Uniform & & & & & & & \\
GPAS (per step) & & & & & & & \\
GPAS (cost-aware) & & & & & & & \\
Unpreconditioned allocation & & & & & & & \\
D$^3$ signal, fixed weights & & & & & & & \\
\midrule
D$^3$-MOPD scheduling & & & & & & & --- \\
Open-MOPD ($K=1$) & & & & & & & --- \\
TA-OPD + Uniform & & & & & & & --- \\
TA-OPD + GPAS & & & & & & & --- \\
\bottomrule
\end{tabular}
\end{table}

\begin{figure}[t]
    \centering
    \missingfigure[0.17\textheight]{Plot common held-out teacher loss
    against optimizer steps and measured GPU hours. Include Uniform,
    the two GPAS modes, and the scheduling comparisons; give the
    recipe-level comparisons in a companion panel.}
    \caption{Learning progress per update and per unit of computation.}
    \label{fig:training_efficiency}
\end{figure}

Table~\ref{tab:main_results} and Figure~\ref{fig:training_efficiency}
evaluate whether the allocation translates into useful training progress.
\missing{Report measured efficiency and every per-domain score change
relative to Uniform. Describe the observed tradeoff between the two
GPAS modes and the recipe-level comparisons.}

\subsection{Where does the computation go?}

\begin{figure}[t]
    \centering
    \missingfigure[0.18\textheight]{For each task, align smoothed gradient
    noise, allocated micro-batches, teacher loss and its rate of change,
    and response length over training. Show preconditioned and
    unpreconditioned noise with clear, separate scales.}
    \caption{Allocation dynamics from the ordinary training trace.
    The panels connect changes in computation to changes in the
    task-gradient estimates and learning progress.}
    \label{fig:allocation_dynamics}
\end{figure}

Figure~\ref{fig:allocation_dynamics} follows the decisions behind the
learning curves. The comparison with unpreconditioned allocation
isolates the optimizer's coordinate scaling; the fixed-weight
D$^3$ signal compares noise with remaining gap and learning velocity.
\missing{Describe the main allocation shifts and their timing.
Explain where the signals agree or disagree using the task traces,
and relate those decisions to the capability and efficiency results.}

\subsection{Cost-aware allocation and reuse with a second dense loss}

The cost-aware run uses the observed time of each task alongside its
gradient variability. We report the fitted step-time components,
observed step times, and the time spent collecting GPAS statistics.
\missing{Report the measured cost differences and the resulting
allocation and GPU-hour differences between the two GPAS modes.}

The TA-OPD pair reuses the same scheduler with tail-aware supervision.
\missing{Report the capability and learning-curve differences between
TA-OPD + Uniform and TA-OPD + GPAS. Compare the resulting allocation
patterns with those under the default teacher top-$64$ loss.}
